% =====================================================================
%  CHAPTER 17: 社区卫生服务
% =====================================================================
\chapter{社区卫生服务}
\setcounter{chapter}{17}
\chapterintro{
掌握社区卫生服务的特点和管理（\textbf{教学难点}）；熟悉社区卫生服务的概念、对象、内容及全科医学的相关概念；了解社区卫生服务的发展概况、方式及绩效评价。
}

% ----- 17.1 社区与社区卫生服务 -----
\section{社区与社区卫生服务概述}

\subsection{社区的定义、类型与要素}

\begin{defbox}
\textbf{社区（Community）}：由一定数量的人群组成，他们居住在共同的地理环境之中，具有共同的文化、共同的生活方式和共同的利益，拥有共同的社会组织和社会服务体系。

\textbf{最早的社区定义}：1887年德国社会学家Ferdinand Tönnies提出"Gemeinschaft"（社区/共同体）概念，区别于"Gesellschaft"（社会）。

\textbf{中国社区的概念}：1933年费孝通等学者将英文"Community"译为"社区"。

\textbf{WHO的社区定义：}一个代表性的社区，其人口数大约在10万-30万之间，面积在0.5万-5万平方公里。
\end{defbox}

\begin{tcolorbox}[colback=light, colframe=main, boxrule=0.5pt, arc=4pt, title=\textbf{社区的要素与类型}, breakable]
\textbf{构成社区的五大要素：}
\begin{enumerate}[leftmargin=*]
    \item \imp{一定数量的人群}——社区的主体
    \item \imp{一定的地域范围}——社区的空间载体
    \item \imp{共同的文化和生活方式}——社区的纽带
    \item \imp{共同的利益和认同感}——社区的凝聚力
    \item \imp{一定的社会组织和服务体系}——社区的运作机制
\end{enumerate}

\textbf{社区类型：}
\begin{itemize}[leftmargin=*]
    \item 按人口分布：城市社区、农村社区、城乡结合部社区
    \item 按功能特征：居住型社区、工作型社区、混合型社区
    \item 按形成方式：自然形成社区、规划设置社区（如街道办事处辖区、居委会辖区）
\end{itemize}
\end{tcolorbox}

\subsection{社区卫生服务的概念与对象}

\begin{defbox}
\textbf{社区卫生服务（Community Health Service, CHS）}：在政府领导、社区参与、上级卫生机构指导下，以\keyword{基层卫生机构为主体}，全科医师为骨干，合理使用社区资源和适宜技术，以人的健康为中心、家庭为单位、社区为范围、需求为导向，以妇女、儿童、老年人、慢性病人、残疾人、贫困居民等为重点服务对象，以解决社区主要卫生问题、满足基本卫生服务需求为目的，融\imp{预防、医疗、保健、康复、健康教育、计划生育技术服务}等于一体的有效、经济、方便、综合、连续的基层卫生服务。
\end{defbox}

\begin{keybox}
\textbf{社区卫生服务的对象：}
\begin{enumerate}[leftmargin=*]
    \item \imp{健康人群}——提供健康促进和疾病预防服务
    \item \imp{亚健康人群}——早期发现健康问题，进行健康管理
    \item \imp{高危人群}——重点监测和干预健康危险因素
    \item \imp{重点保健人群}——妇女、儿童、老年人
    \item \imp{患者}——常见病、多发病的诊治和慢性病管理
\end{enumerate}
\end{keybox}

\subsection{全科医学}

\begin{defbox}
\textbf{全科医学/家庭医学（General Practice / Family Medicine）}：整合临床医学、预防医学、康复医学及人文社会科学相关内容于一体的综合性医学专业学科，是服务于社区、以家庭为单位、以个人为中心的\keyword{连续性、综合性、协调性}的基层医疗保健专业。

\textbf{全科医师（General Practitioner, GP）}：经过全科医学专业培训的医生，是社区卫生服务的骨干力量，提供第一线的全面、连续的卫生保健服务。
\end{defbox}

% ----- 17.2 社区卫生服务的特点 -----
\section{社区卫生服务的特点与内容}

\subsection{社区卫生服务的特点（教学难点*）}

\begin{keybox}
\imp{（教学难点*）}\textbf{社区卫生服务的六大特点：}

\begin{enumerate}[leftmargin=*]
    \item \imp{以健康为中心（Health-centred）}——以保护和促进人群健康为目标，而非仅仅"治病"。强调"以人的健康为中心"，关注全人群和全生命周期健康

    \item \imp{以需求为导向（Demand-oriented）}——根据社区居民的实际健康需求和主要健康问题来规划和组织服务，服务内容和方式与社区需求相匹配

    \item \imp{以家庭为单位（Family as a Unit）}——将家庭作为服务的基本单元，关注家庭结构、功能对成员健康的影响。家庭是健康行为养成和社会支持的核心场域：
    \begin{itemize}
        \item 家庭影响健康行为的形成和维持
        \item 家庭关系提供情感支持和疾病照护
        \item 家庭功能障碍可成为疾病诱因
        \item "一人患病，全家受累"——家庭疾病负担传递
    \end{itemize}

    \item \imp{以社区为范围（Community-based）}——服务涵盖整个社区，关注社区整体健康状况和社区环境对健康的影响

    \item \imp{预防为主（Prevention-focused）}——将预防贯穿于医疗卫生服务的始终。社区是三级预防的最佳场所：一级预防（健康教育、预防接种）、二级预防（筛查、早期诊断）、三级预防（慢性病管理、康复）

    \item \imp{综合性、连续性和协调性服务}——六位一体综合服务（预防、医疗、保健、康复、健康教育、计划生育）。连续性——从围生期到临终的生命全程服务；首次接触（First Contact）——社区卫生服务是居民进入卫生系统的第一站，解
\end{enumerate}
\end{keybox}

\subsection{社区卫生服务的内容}

\begin{enumerate}[leftmargin=*]
    \item \imp{预防服务}：健康教育、计划免疫、传染病防控、慢性病预防与控制
    \item \imp{医疗服务}：常见病和多发病的诊疗、院前急救、家庭病床
    \item \imp{保健服务}：妇女保健、儿童保健、老年保健、重点人群健康管理
    \item \imp{康复服务}：社区康复、慢性病康复、残疾人士社区康复
    \item \imp{健康教育}：健康知识普及、健康行为指导、健康技能培训
    \item \imp{计划生育技术服务}：避孕节育指导、优生优育咨询
\end{enumerate}

% ----- 17.3 社区卫生服务发展 -----
\section{社区卫生服务发展概况}

\subsection{国际发展}

\begin{itemize}[leftmargin=*]
    \item 1945年：英国工党政府开始建立NHS（国家卫生服务），确立全科医师"守门人"制度
    \item 20世纪50-60年代：日本、澳大利亚等国先后建立社区/家庭医学体系
    \item 1978年：Alma-Ata宣言推动全球社区卫生服务和初级卫生保健发展
\end{itemize}

\subsection{中国发展}

\begin{enumerate}[leftmargin=*]
    \item 1997年：《中共中央、国务院关于卫生改革与发展的决定》明确发展社区卫生服务
    \item 1999年：十部委发布《关于发展城市社区卫生服务的若干意见》
    \item 2002年：十一部委发布《关于加快发展城市社区卫生服务的意见》
    \item 2006年：《国务院关于发展城市社区卫生服务的指导意见》出台
    \item 2009年：新医改将社区卫生服务作为"保基本、强基层、建机制"的重点
    \item 2011年：全面推行家庭医生签约服务
    \item 2015年：《关于推进分级诊疗制度建设的指导意见》强化社区首诊
    \item 2016年：《"健康中国2030"规划纲要》将社区卫生服务作为基层医疗核心
    \item 2018年：《关于加强公立医院党的建设工作的意见》推进医联体建设
\end{enumerate}

% ----- 17.4 社区卫生服务管理 -----
\section{社区卫生服务的管理（教学难点*）}

\begin{keybox}
\imp{（教学难点*）}\textbf{社区卫生服务管理的核心要点：}

\textbf{1. 组织管理：}
\begin{itemize}[leftmargin=*]
    \item 社区卫生服务中心（站）的设置标准（服务人口：城市3-10万人/中心，农村乡镇）
    \item 运行机制：政府主导、公益性原则
\end{itemize}

\textbf{2. 人员管理：}
\begin{itemize}[leftmargin=*]
    \item 全科医师（GP）为核心的团队建设
    \item 家庭医生签约服务团队（GP+护士+公卫医师等）
    \item 继续教育和能力提升
\end{itemize}

\textbf{3. 服务质量管理：}
\begin{itemize}[leftmargin=*]
    \item 临床路径和诊疗规范
    \item 双向转诊制度
    \item 绩效考核和质量评价
\end{itemize}

\textbf{4. 经费与医保管理：}
\begin{itemize}[leftmargin=*]
    \item 财政投入保障（政府购买服务、基本公共卫生服务均等化资金）
    \item 医保支付方式改革（按人头付费、总额预付等）
\end{itemize}
\end{keybox}

\subsection{服务方式与绩效评价}

\begin{itemize}[leftmargin=*]
    \item \imp{服务方式}：门诊服务、上门服务（家庭病床、出诊）、社区巡诊、电话咨询服务、网络健康管理、家庭医生签约服务
    \item \imp{绩效评价}：服务覆盖率（健康档案建档率、计划免疫接种率等）、服务质量（慢性病管理率、血压/血糖控制率等）、居民满意度、服务效率（门诊人次、住院床日等）、费用控制效果
\end{itemize}

\subsection{社区卫生服务的意义}

\begin{enumerate}[leftmargin=*]
    \item \imp{体现社会公平}——使全体居民特别是弱势群体享受到基本卫生服务
    \item \imp{合理利用卫生资源}——常见病多发病在基层解决，缓解大医院"看病难"
    \item \imp{提高卫生服务效率}——通过分级诊疗和双向转诊优化资源配置
    \item \imp{降低医疗费用}——社区首诊和健康管理减少不必要的高成本服务
    \item \imp{实现医学模式转变}——社区卫生服务是生物-心理-社会医学模式的最佳实践场所
\end{enumerate}

% ----- 复习题 -----
\section{复习题}

\subsection*{一、选择题}
\begin{enumerate}[leftmargin=*, itemsep=3pt]
    \item 社区卫生服务的核心骨干是：
    \begin{enumerate}
        \item[(A)] 专科医师
        \item[(B)] 全科医师
        \item[(C)] 公共卫生医师
        \item[(D)] 护士
    \end{enumerate}
    \item 社区卫生服务的特点不包括：
    \begin{enumerate}
        \item[(A)] 以健康为中心
        \item[(B)] 以家庭为单位
        \item[(C)] 以专科医疗为重点
        \item[(D)] 预防为主
    \end{enumerate}
    \item 社区卫生服务的"六位一体"不包括：
    \begin{enumerate}
        \item[(A)] 预防
        \item[(B)] 医疗
        \item[(C)] 器官移植
        \item[(D)] 健康教育
    \end{enumerate}
    \item 将"Community"译为"社区"的中国学者是：
    \begin{enumerate}
        \item[(A)] 陈志潜
        \item[(B)] 费孝通
        \item[(C)] 吴阶平
        \item[(D)] 兰安生
    \end{enumerate}
\end{enumerate}

\subsection*{二、名词解释}
\begin{enumerate}[leftmargin=*, itemsep=3pt]
    \item \textbf{社区卫生服务（CHS）}
    \item \textbf{全科医学（General Practice）}
    \item \textbf{社区（Community）}
\end{enumerate}

\subsection*{三、简答题}
\begin{enumerate}[leftmargin=*, itemsep=3pt]
    \item 试述社区卫生服务的六大特点。
    \item 简述社区卫生服务的"六位一体"内容。
    \item 简述社区卫生服务的重要意义。
    \item 简述我国社区卫生服务的发展历程。
\end{enumerate}

\subsection*{参考答案要点}

\subsubsection*{一、选择题}
1. (B)；2. (C)；3. (C)；4. (B)

\subsubsection*{二、名词解释}
\begin{enumerate}[leftmargin=*]
    \item \textbf{CHS}：以基层卫生机构为主体、全科医师为骨干，融预防、医疗、保健、康复、健康教育、计划生育技术服务于一体，以健康为中心、家庭为单位、社区为范围的基层卫生服务。
    \item \textbf{全科医学}：整合临床、预防、康复及人文社会科学的综合性医学专业学科，提供连续性、综合性、协调性的基层医疗保健。
    \item \textbf{社区}：由一定数量人群组成，居住在共同地域，具有共同文化、生活方式和利益，拥有共同社会组织和服务体系的社会单位。
\end{enumerate}

\subsubsection*{三、简答题}
\begin{enumerate}[leftmargin=*]
    \item 六大特点：①以健康为中心（非仅"治病"）；②以需求为导向（服务与需求匹配）；③以家庭为单位（关注家庭对健康的影响）；④以社区为范围（关注社区整体健康）；⑤预防为主（三级预防的理想场所）；⑥综合性、连续性和协调性（六位一体+生命全程+首诊制）。
    \item 六位一体：预防（健康教育+免疫接种+传染病防控）、医疗（常见病诊疗+急救）、保健（妇幼+老年+重点人群）、康复（社区康复+慢性病康复）、健康教育（知识普及+行为指导）、计划生育技术服务（避孕节育+优生优育）。
    \item 五大意义：体现社会公平（弱势群体可及性）、合理利用资源（基层首诊缓解大医院压力）、提高效率（分级诊疗+双向转诊）、降低费用（社区首诊减少高成本服务）、实现医学模式转变（BPS模式最佳实践场所）。
    \item 重要里程碑：1997年中央决定发展CHS；1999年十部委意见；2006年国务院指导意见；2009年新医改"保基本强基层"; 2011年家庭医生签约；2015年分级诊疗制度；2016年健康中国2030。
\end{enumerate}

\newpage
